\section{Vorbereitung der Autoren}
Die letzten Sonntage habe ich erzählt, dass die Bibel Gottes Wort ist, dass diese fehlerlos ist, dass diese für alles genügt, also allgenügsam ist.

Aber es nützt ja nichts, wenn ich es euch von hier vorne erzähle. Wer bin ich, dass ich das einfach so behaupten kann? Darum wollen wir näher untersuchen, ob meine Behauptungen stimmen können.

Ein berechtigter Einwand: \enquote{Die Bibel wurde doch von Menschen geschrieben. Menschen sind Sünder und machen Fehler. Ausserdem sind jetzt 4000 Jahre vergangen und viele haben an den Texten herumgespielt? Wie kann der noch stimmen?}

Richtig das stimmt. Alle der 40 Autoren, die an diesem Buch mitgeschrieben haben, waren \enquote{normale} Menschen wie du und ich. Sie waren Sünden wie du und ich und es waren Menschen, die genauso das Blut Christi brauchen für ihre Errettung wie du und ich. Und genau das ist ein Aspekt, der mich überzeugt, dass die Autoren von Gott geleitet wurden. Kennt ihr noch den Heilsplan Gottes? Wir haben den letztes Jahr hier uns zusammen angeschaut. 

\textit{Ich bin übrigens dabei, eine Zusammenstellung der Vorträge zu machen und werde die dann an euch abgeben. Frage ist nur wann. } 

Kein Mensch kann über 1700 Jahre ein Buch zusammen stellen, das so in sich  stimmig ist, aufbaut und ergänzt. Und die in diesem Buch enthalten Prophetien habe ich noch nicht mal erwähnt. 

Es waren keine Übermenschen am Werk. Es gab ein paar Autoren, die eine spezielle Ausbildung hatten, und welche die einen hohen Posten eingenommen haben. Wiederum andere die Bauern waren oder Fischer. Aber auch das sind keine Zufälle. Nehmen wir als Beispiel Mose:
\begin{enumerate}
    \item Genau zur richtigen Zeit als Jude geboren. \biblezit{IMos}{7:1}
    \item Genau zu der Zeit gab es eine Verfolgung in Ägypten. Darum wurde er im Nil ausgesetzt.
    \item Genau zu der Zeit war die Tochter des Pharao am Nil baden und fand den Säugling
    \item Er wurde am Hof des Pharao ausgebildet 40 Jahre.
    \item Er wurde in Wüste 40 Jahre ausgebildet. 
    \item Und dann kam Gott in einem brennenden Dornbusch.
\end{enumerate}
Gott war zur Ausführung seiner Absicht, nie von den Schreibern abhängig. So sehen wir zum Beispiel, dass Gott Mose nicht in das verheissene Land einziehen liess. Es ging trotzdem weiter. Einer Führer und Schreiber kam der von Mose ausgebildet wurde.
Die ersten 5 Bücher der Bibel sind gleichzeitig das Werk Gottes und Moses. Wissenschaftliche Untersuchungen und Textvergleiche haben gezeigt, dass Mose viele ägyptische Redeweisen und Begriffen in seine Schriften benutzt hat.

Im neuen Testament haben vier verschiedene Personen die die Evangelien geschrieben haben. Drei synoptische Evangelium und das Johannes Evangelium. Synoptisch heisst, \enquote{zusammenschauend}, (griechisch synoptikos = zusammen-sehend).

Seit 2000 Jahren wird gestritten, wer von wem abgeschrieben hat und wer der erste war. Es wird versucht sie gegeneinander auszuspielen.

Es gab hunderttausende, wenn nicht Millionen von Wissenschaftler, Theologen, Religionsstifter und sonstige, die versuchen Fehler in der Bibel zu finden und so die Richtigkeit und Glaubwürdigkeit der Evangelien zu widerlegen. Na ja und das Resultat? Der Text ist immer noch der gleiche, den schon die Kirchenväter im 2. Jahrhundert nach Christus lasen.

Das alles waren keine zufälligen Schreiberlinge. So lesen wir in \biblezit{Jer}{1:5} 
\begin{bibelbox}{SCHL}{Jer}{1:5}
Ehe ich dich im Mutterleib bildete, habe ich dich ersehen, und bevor du aus dem Mutterschoss hervorkamst, habe ich dich geheiligt; zum Propheten für die Völker habe ich dich bestimmt!
\end{bibelbox}
und \biblezit{Gal}{1:15}
\begin{bibelbox}{SCHL}{Gal}{1:15}
Als es aber Gott, der mich vom Mutterleib an ausgesondert und durch seine Gnade berufen hat,
\end{bibelbox}
, dass Paulus und Jeremia schon vor ihrer Geburt von Gott erwählt wurden. So auch alle andern Autoren. Auch Jesus sagt in \biblezit{Joh}{15:16}
\begin{bibelbox}{SCHL}{Joh}{15:16}
Nicht ihr habt mich erwählt, sondern ich habe euch dazu bestimmt...
\end{bibelbox}

John McArthur schreibt dazu:
\begin{quote}    
    Alle diese Faktoren wirken zusammen, um diese Männer zu den Werkzeugen zu formen, zu denen Gott sie bestimmt hat, um göttlich autoritative Schriften hervorzubringen.
    
    Alle diese einzigartigen sozialen, kulturellen, historischen, emotionellen, erfahrungsmässigen, Bildungsmässigen und praktischen Eigenheiten spiegeln sich in der Sprache und dem Stil des Werkes jedes menschlichen Autors wider.
\end{quote}

\begin{otsblock}[Fazit]
\begin{enumerate}
    \item Alle Autoren waren Sünder und Menschen wie du und ich.
    \item Alle Autoren hatten einen eigenen geschichtlichen Hintergrund.
    \item Diese Eigenarten sind in die Texte mit eingeflossen.
    \item Gott hat die Schreiber schon vor ihrer Geburt auserwählt.
    \item Durch Gottes Regie entstand ein in sich aufbauendes stimmiges Werk, dass noch immer so aktuell ist wie zu der Zeit als es geschrieben wurde.
\end{enumerate}
Nächsten Sonntag schauen wir uns die alttestamentlichen Beweise an. Das ganze muss jetzt noch bewiesen werden.
\end{otsblock}